%%% Zahlen und Rechenarten
\begin{lesen}
\lies{$12 + 30 = 42$}{The sum of twelve and thirty is forty-two.}{Die Summe von zwölf und dreißig ist zweiundvierzig.}
\lies{$3 \mid 12$}{Three is a divisor (a factor) of twelve; twelve is a multiple of three; twelve is divisible by three.}{Drei ist ein Teiler von zwölf; zwölf ist ein Vielfaches von drei; zwölf ist durch drei teilbar.}
\lies{$\operatorname{ggT}(12;\,18) = 6$}{The greatest common divisor of twelve and eighteen is six.}{Der größte gemeinsame Teiler von zwölf und achtzehn ist sechs.}
\lies{$17$}{Seventeen is a prime number: it has exactly two divisors, one and itself.}{Siebzehn ist eine Primzahl: Sie hat genau zwei Teiler, die Eins und sich selbst.}
\lies{$3{,}14159 \approx 3{,}14$}{Round the number to two decimal places.}{Runde die Zahl auf zwei Nachkommastellen.}
\lies{$|-7| = 7$}{The absolute value of minus seven is seven.}{Der Betrag von minus sieben ist sieben.}
\lies{$25\,\%$ von $80$ sind $20$}{Twenty-five per cent of eighty is twenty.}{Fünfundzwanzig Prozent von achtzig sind zwanzig.}
\end{lesen}

%%% Algebra, Gleichungen und Mengen
\begin{lesen}
\lies{$3(x + 2) = 3x + 6$}{Expand the brackets.}{Multipliziere die Klammer aus.}
\lies{$6x + 9 = 3(2x + 3)$}{Factorise the expression: take out the common factor three.}{Klammere aus: Der gemeinsame Faktor ist drei.}
\lies{$(x + 3)^2 = x^2 + 6x + 9$}{$x$ plus three, all squared, equals $x$ squared plus six $x$ plus nine.}{$x$ plus drei in Klammern zum Quadrat ist gleich $x$ Quadrat plus sechs $x$ plus neun.}
\lies{$2x + 3 = 11$ \newline $x = 4$}{Solve the equation for $x$. The solution is $x$ equals four.}{Löse die Gleichung nach $x$ auf. Die Lösung ist $x$ gleich vier.}
\lies{$x + y = 5$ \newline $x - y = 1$}{Solve the system of equations: $x$ equals three and $y$ equals two.}{Löse das Gleichungssystem: $x$ gleich drei und $y$ gleich zwei.}
\lies{$A = \{1;\,2;\,3\},\ B = \{3;\,4\}$ \newline $A \cap B = \{3\}$}{The intersection of $A$ and $B$ contains only the element three.}{Die Schnittmenge von $A$ und $B$ enthält nur das Element drei.}
\end{lesen}

%%% Funktionen
\begin{lesen}
\lies{$f(x) = 2x - 4$}{$f$ is a linear function with the gradient two and the $y$-intercept minus four. Its zero is $x$ equals two.}{$f$ ist eine lineare Funktion mit der Steigung zwei und dem $y$-Achsenabschnitt minus vier. Ihre Nullstelle ist $x$ gleich zwei.}
\lies{$f(x) = (x - 1)^2 - 4$}{The graph is a parabola with the vertex $S(1 \mid -4)$. The zeros are minus one and three.}{Der Graph ist eine Parabel mit dem Scheitelpunkt $S(1 \mid -4)$. Die Nullstellen sind minus eins und drei.}
\lies{$f(x) = x^2$}{The graph is symmetric about the $y$-axis. The function is decreasing for $-\infty < x \le 0$ and increasing for $0 \le x < \infty$.}{Der Graph ist achsensymmetrisch zur $y$-Achse. Die Funktion ist fallend für $-\infty < x \le 0$ und steigend für $0 \le x < \infty$.}
\lies{$f(x) = 2^x$}{The graph of this exponential function passes through the point $(0 \mid 1)$; the $x$-axis is an asymptote.}{Der Graph dieser Exponentialfunktion verläuft durch den Punkt $(0 \mid 1)$; die $x$-Achse ist Asymptote.}
\lies{$D_f = \R$ \newline $W_f = \{y \in \R \mid y \ge -4\}$}{The domain is the set of all real numbers; the range is the set of all real numbers greater than or equal to minus four.}{Der Definitionsbereich ist die Menge aller reellen Zahlen; der Wertebereich ist die Menge aller reellen Zahlen, die größer oder gleich minus vier sind.}
\end{lesen}

%%% Geometrie: Grundbegriffe, Winkel, Abbildungen
\begin{lesen}
\lies{$\overline{AB} = 5\,\mathrm{cm}$}{The line segment $A$ $B$ is five centimetres long.}{Die Strecke $A$ $B$ ist fünf Zentimeter lang.}
\lies{$g \parallel h$}{Draw a line $h$ through the point $P$ parallel to $g$.}{Zeichne durch den Punkt $P$ die Parallele $h$ zu $g$.}
\lies{$\alpha = 35\degree$}{Alpha is an acute angle of thirty-five degrees.}{Alpha ist ein spitzer Winkel von fünfunddreißig Grad.}
\lies{$\beta = 180\degree - 70\degree = 110\degree$}{Adjacent angles on a straight line add up to one hundred and eighty degrees.}{Nebenwinkel ergänzen sich zu hundertachtzig Grad.}
\lies{\skSpiegelung}{Reflect the triangle in the line $g$.}{Spiegele das Dreieck an der Geraden $g$.}
\lies{\skDrehung}{Rotate the triangle through ninety degrees about the point $Z$.}{Drehe das Dreieck um neunzig Grad um den Punkt $Z$.}
\lies{\skMittelsenkrechte}{Construct the perpendicular bisector of the line segment $A$ $B$.}{Konstruiere die Mittelsenkrechte der Strecke $A$ $B$.}
\end{lesen}

%%% Ebene Figuren
\begin{lesen}
\lies{$a = 6\,\mathrm{cm},\ b = 4\,\mathrm{cm}$ \newline $A = 24\,\mathrm{cm}^2$}{The area of the rectangle is six centimetres times four centimetres, that is twenty-four square centimetres.}{Der Flächeninhalt des Rechtecks ist sechs Zentimeter mal vier Zentimeter, also vierundzwanzig Quadratzentimeter.}
\lies{$3^2 + 4^2 = 5^2$}{A triangle with the sides three, four and five is right-angled.}{Ein Dreieck mit den Seitenlängen drei, vier und fünf ist rechtwinklig.}
\lies{\skGleichschenklig}{An isosceles triangle has two equal sides and two equal base angles.}{Ein gleichschenkliges Dreieck hat zwei gleich lange Seiten und zwei gleich große Basiswinkel.}
\lies{\skDiagonale}{The diagonals of a rectangle are equal in length and bisect each other.}{Die Diagonalen eines Rechtecks sind gleich lang und halbieren einander.}
\lies{$d = 10\,\mathrm{cm}$ \newline $u = \pi \cdot d \approx 31{,}4\,\mathrm{cm}$}{The circumference of a circle with a diameter of ten centimetres is approximately thirty-one point four centimetres.}{Der Umfang eines Kreises mit zehn Zentimetern Durchmesser beträgt ungefähr einunddreißig Komma vier Zentimeter.}
\end{lesen}

%%% Körper
\begin{lesen}
\lies{$a = 3\,\mathrm{cm}$ \newline $V = a^3 = 27\,\mathrm{cm}^3$}{The volume of a cube with an edge length of three centimetres is twenty-seven cubic centimetres.}{Das Volumen eines Würfels mit drei Zentimetern Kantenlänge beträgt siebenundzwanzig Kubikzentimeter.}
\lies{$O = 6a^2 = 54\,\mathrm{cm}^2$}{Its surface area is fifty-four square centimetres.}{Sein Oberflächeninhalt beträgt vierundfünfzig Quadratzentimeter.}
\lies{\skNetz}{The net of a cube consists of six squares.}{Das Netz eines Würfels besteht aus sechs Quadraten.}
\lies{\skPyramide}{A square-based pyramid has five vertices, eight edges and five faces.}{Eine quadratische Pyramide hat fünf Ecken, acht Kanten und fünf Flächen.}
\lies{$r = 2\,\mathrm{cm},\ h = 5\,\mathrm{cm}$ \newline $V = \pi r^2 h \approx 62{,}8\,\mathrm{cm}^3$}{The volume of the cylinder is twenty pi, approximately sixty-two point eight cubic centimetres.}{Das Volumen des Zylinders ist zwanzig pi, ungefähr zweiundsechzig Komma acht Kubikzentimeter.}
\end{lesen}

%%% Analysis
\begin{lesen}
\lies{$f(x) = x^3 - 3x$ \newline $f'(x) = 3x^2 - 3$}{Differentiate $f$: the derivative is three $x$ squared minus three.}{Leite $f$ ab: Die Ableitung ist drei $x$ Quadrat minus drei.}
\lies{$f'(x) = 0 \ \Rightarrow\ x = \pm 1$}{The derivative is zero at $x$ equals plus or minus one.}{Die Ableitung ist null für $x$ gleich plus oder minus eins.}
\lies{$H(-1 \mid 2),\ T(1 \mid -2)$ \newline $W(0 \mid 0)$}{The graph has a maximum at minus one two, a minimum at one minus two and a point of inflection at the origin.}{Der Graph hat den Hochpunkt $H$, den Tiefpunkt $T$ und den Wendepunkt $W$ im Ursprung.}
\lies{$F(x) = x^2$ \newline $f(x) = 2x$}{Capital $F$ is an antiderivative of $f$.}{Groß $F$ ist eine Stammfunktion von $f$.}
\lies{$\displaystyle\int_0^3 2x \dx = \big[x^2\big]_0^3 = 9$}{The integral from nought to three of two $x$ d $x$ equals nine.}{Das Integral von null bis drei über zwei $x$ d $x$ ist neun.}
\lies{$\limn \dfrac{2n + 1}{n} = 2$}{The sequence converges; its limit is two.}{Die Folge konvergiert; ihr Grenzwert ist zwei.}
\end{lesen}

%%% Vektoren und analytische Geometrie
\begin{lesen}
\lies{$\vec{a} = \vektor{3}{4}{0}$ \newline $|\vec{a}\,| = \sqrt{9 + 16 + 0} = 5$}{The magnitude of vector $a$ is five.}{Der Betrag des Vektors $a$ ist fünf.}
\lies{$A(1 \mid 2 \mid 3),\ B(4 \mid 6 \mid 3)$ \newline $\overrightarrow{AB} = \vektor{3}{4}{0}$}{The vector from $A$ to $B$ is three, four, nought.}{Der Vektor von $A$ nach $B$ ist drei, vier, null.}
\lies{$g\colon\ \vec{x} = \vektor{1}{2}{3} + t \cdot \vektor{3}{4}{0}$}{The equation of the line $g$: a position vector plus $t$ times a direction vector.}{Die Gleichung der Geraden $g$: Stützvektor plus $t$ mal Richtungsvektor.}
\lies{$\vec{a} \cdot \vec{b} = 0$}{The vectors $a$ and $b$ are perpendicular (orthogonal).}{Die Vektoren $a$ und $b$ sind zueinander senkrecht (orthogonal).}
\lies{$\vec{b} = 2 \cdot \vec{a}$}{The vectors $a$ and $b$ are linearly dependent (parallel).}{Die Vektoren $a$ und $b$ sind linear abhängig (parallel).}
\end{lesen}

%%% Stochastik
\begin{lesen}
\lies{$P(\text{„6“}) = \dfrac16$}{The probability of rolling a six with a fair die is one sixth.}{Die Wahrscheinlichkeit, mit einem fairen Würfel eine Sechs zu würfeln, beträgt ein Sechstel.}
\lies{$P(\text{„zweimal Kopf“}) = \dfrac12 \cdot \dfrac12 = \dfrac14$}{A fair coin is tossed twice. The probability of getting two heads is one quarter.}{Eine faire Münze wird zweimal geworfen. Die Wahrscheinlichkeit für zweimal Kopf beträgt ein Viertel.}
\lies{$P(\overline{A}) = 1 - P(A)$}{The probability of the complementary event equals one minus the probability of $A$.}{Die Wahrscheinlichkeit des Gegenereignisses ist eins minus die Wahrscheinlichkeit von $A$.}
\lies{$2;\ 4;\ 6;\ 8$ \newline $\bar{x} = 5$}{The mean of the four values is five.}{Der Mittelwert der vier Werte ist fünf.}
\lies{$\E(X) = 3{,}5$}{The expected value of the number shown by a fair die is three point five.}{Der Erwartungswert der Augenzahl eines fairen Würfels ist drei Komma fünf.}
\end{lesen}

%%% Logik und Beweis
\begin{lesen}
\lies{$x > 2 \ \Rightarrow\ x^2 > 4$}{If $x$ is greater than two, then $x$ squared is greater than four.}{Wenn $x$ größer als zwei ist, dann ist $x$ Quadrat größer als vier.}
\lies{$n$ gerade $\Leftrightarrow$ $n^2$ gerade}{$n$ is even if and only if $n$ squared is even.}{$n$ ist genau dann gerade, wenn $n$ Quadrat gerade ist.}
\lies{$\forall\, x \in \R\colon\ x^2 \ge 0$}{For all real numbers $x$, $x$ squared is greater than or equal to zero.}{Für alle reellen Zahlen $x$ gilt: $x$ Quadrat ist größer oder gleich null.}
\lies{$\exists\, x \in \R\colon\ x^2 = 2$}{There exists a real number $x$ such that $x$ squared equals two.}{Es gibt eine reelle Zahl $x$, für die $x$ Quadrat gleich zwei ist.}
\lies{Annahme}{Assume that the square root of two is rational. \ldots\ This is a contradiction. (proof by contradiction)}{Angenommen, die Wurzel aus zwei ist rational. \ldots\ Das ist ein Widerspruch. (Beweis durch Widerspruch)}
\end{lesen}

%%% Zahlen, Zeichen und Schreibweisen
\begin{lesen}
\lies{$7{,}25$}{seven point two five}{sieben Komma zwei fünf}
\lies{$0{,}05$}{nought point nought five}{null Komma null fünf}
\lies{$1\,250\,000$ \newline \footnotesize engl. 1,250,000}{one million two hundred and fifty thousand}{eine Million zweihundertfünfzigtausend}
\end{lesen}

%%% Relationszeichen
\begin{lesen}
\lies{$5 < 8$}{five is less than eight}{fünf ist kleiner als acht}
\lies{$x \ge 0$}{$x$ is greater than or equal to zero: $x$ is non-negative}{$x$ ist größer oder gleich null: $x$ ist nicht negativ}
\lies{$7 \neq 9$}{seven is not equal to nine}{sieben ist ungleich neun}
\lies{$\sqrt{2} \approx 1{,}41$}{the square root of two is approximately one point four one}{die Wurzel aus zwei ist ungefähr eins Komma vier eins}
\lies{$-2 < x < 2$}{$x$ lies between minus two and two}{$x$ liegt zwischen minus zwei und zwei}
\end{lesen}

%%% Rechnen mit Brüchen
\begin{lesen}
\lies{$\dfrac56 - \dfrac14 = \dfrac{10}{12} - \dfrac{3}{12} = \dfrac{7}{12}$}{five sixths minus one quarter equals ten twelfths minus three twelfths, which is seven twelfths}{fünf Sechstel minus ein Viertel ist gleich zehn Zwölftel minus drei Zwölftel, also sieben Zwölftel}
\lies{$\dfrac45 \cdot 10 = 8$}{four fifths of ten is eight}{vier Fünftel von zehn sind acht}
\lies{$2\tfrac12 = \dfrac52$}{two and a half equals five halves}{zweieinhalb ist gleich fünf Halbe}
\lies{$\dfrac38 = 0{,}375 = 37{,}5\,\%$}{three eighths equals nought point three seven five or thirty-seven point five per cent}{drei Achtel ist gleich null Komma drei sieben fünf oder siebenunddreißig Komma fünf Prozent}
\end{lesen}

%%% Potenzen, Wurzeln und Logarithmen
\begin{lesen}
\lies{$2^{10} = 1024$}{two to the power of ten equals one thousand and twenty-four}{zwei hoch zehn ist eintausendvierundzwanzig}
\lies{$a^m \cdot a^n = a^{m+n}$}{$a$ to the $m$ times $a$ to the $n$ equals $a$ to the $m$ plus $n$}{$a$ hoch $m$ mal $a$ hoch $n$ ist gleich $a$ hoch $m$ plus $n$}
\lies{$9^{\frac12} = \sqrt{9} = 3$}{nine to the power of one half equals the square root of nine, which is three}{neun hoch ein Halb ist gleich der Wurzel aus neun, also drei}
\lies{$\sqrt[3]{27} = 3$}{the cube root of twenty-seven is three}{die dritte Wurzel aus siebenundzwanzig ist drei}
\end{lesen}

%%% Zahlbereiche
\begin{lesen}
\lies{$-3 \in \Z,\ -3 \notin \N$}{minus three is an integer, but it is not a natural number}{minus drei ist eine ganze Zahl, aber keine natürliche Zahl}
\lies{$\frac34 \in \Q$}{three quarters is a rational number}{drei Viertel ist eine rationale Zahl}
\lies{$\sqrt{2} \notin \Q$}{the square root of two is irrational}{die Wurzel aus zwei ist irrational}
\lies{$0{,}\overline{3} = \frac13$}{nought point three recurring equals one third}{null Komma Periode drei ist gleich ein Drittel}
\end{lesen}

%%% Mengen
\begin{lesen}
\lies{$A = \{1;\,2;\,3;\,4\}$ \newline $B = \{3;\,4;\,5\}$}{$A$ is the set with the elements one, two, three and four; $B$ is the set with the elements three, four and five}{$A$ ist die Menge mit den Elementen eins, zwei, drei und vier; $B$ ist die Menge mit den Elementen drei, vier und fünf}
\lies{$A \cap B = \{3;\,4\}$}{the intersection of $A$ and $B$ is the set containing three and four}{die Schnittmenge von $A$ und $B$ ist die Menge mit den Elementen drei und vier}
\lies{$A \cup B = \{1;\,2;\,3;\,4;\,5\}$}{the union of $A$ and $B$ contains the numbers from one to five}{die Vereinigungsmenge von $A$ und $B$ enthält die Zahlen von eins bis fünf}
\end{lesen}

%%% Gleichungen und Ungleichungen lösen
\begin{lesen}
\lies{$x^2 - 5x + 6 = 0$ \newline $(x - 2)(x - 3) = 0$}{Factorise the left-hand side.}{Zerlege die linke Seite in Faktoren.}
\lies{$x_1 = 2,\ x_2 = 3$}{A product is zero if one of its factors is zero. The solutions are two and three.}{Ein Produkt ist null, wenn einer der Faktoren null ist. Die Lösungen sind zwei und drei.}
\end{lesen}

%%% Funktionen
\begin{lesen}
\lies{$f(x) = 2x^2 - 3$ \newline $f(2) = 5$}{$f$ of two equals five: the point $P(2 \mid 5)$ lies on the graph of $f$}{$f$ von zwei ist fünf: Der Punkt $P(2 \mid 5)$ liegt auf dem Graphen von $f$}
\lies{$f(x) = \frac12 x + 1$}{a linear function with the gradient one half and the $y$-intercept one}{eine lineare Funktion mit der Steigung ein Halb und dem $y$-Achsenabschnitt eins}
\lies{$f(x) = 0$}{Find the zeros of $f$: for which values of $x$ is $f$ of $x$ equal to zero?}{Bestimme die Nullstellen von $f$: Für welche $x$ ist $f$ von $x$ gleich null?}
\lies{$f(x) = g(x)$}{Find the points of intersection of the two graphs.}{Bestimme die Schnittpunkte der beiden Graphen.}
\end{lesen}

%%% Winkelfunktionen
\begin{lesen}
\lies{$\cos 60\degree = \dfrac12$}{the cosine of sixty degrees is one half}{der Kosinus von sechzig Grad ist ein Halb}
\lies{$\sin x = 0{,}5$ \newline $x_1 = 30\degree,\ x_2 = 150\degree$}{For angles between nought and three hundred and sixty degrees the equation has two solutions: thirty degrees and one hundred and fifty degrees.}{Für Winkel zwischen null und dreihundertsechzig Grad hat die Gleichung zwei Lösungen: dreißig Grad und hundertfünfzig Grad.}
\lies{$\sin(x + 360\degree) = \sin x$}{The sine function is periodic with a period of three hundred and sixty degrees.}{Die Sinusfunktion ist periodisch mit der Periode dreihundertsechzig Grad.}
\lies{$-1 \le \sin x \le 1$}{The values of the sine function lie between minus one and one.}{Die Werte der Sinusfunktion liegen zwischen minus eins und eins.}
\end{lesen}

%%% Intervalle und Grenzwerte
\begin{lesen}
\lies{$x \in [2;\,5]$}{$x$ is an element of the closed interval from two to five}{$x$ ist Element des abgeschlossenen Intervalls von zwei bis fünf}
\lies{$\lim\limits_{x \to \infty} \dfrac{1}{x} = 0$}{the limit of one over $x$ as $x$ tends to infinity is zero}{der Grenzwert von eins durch $x$ für $x$ gegen unendlich ist null}
\lies{$\limn \dfrac{2n + 1}{n} = 2$}{the limit of two $n$ plus one, all over $n$, as $n$ tends to infinity is two}{der Grenzwert von zwei $n$ plus eins, geteilt durch $n$, für $n$ gegen unendlich ist zwei}
\lies{steigend: $2 \le x < \infty$}{the function is increasing for all $x$ greater than or equal to two}{die Funktion ist steigend für alle $x$, die größer oder gleich zwei sind}
\end{lesen}
